\corrigesChapitre{5}{Métaux : réactions et corrosion}
\par\noindent\begin{minipage}{\linewidth}
\begin{corrige}{1}{Corrosion}
\begin{questions}
\item Chimique.
\item Eau et dioxygène.
\end{questions}
\end{corrige}
\end{minipage}
\par\noindent\begin{minipage}{\linewidth}
\begin{corrige}{2}{Un gaz}
\begin{questions}
\item Dihydrogène, $\ce{H2}$.
\item Ion zinc, $\ce{Zn^{2+}}$.
\end{questions}
\pcFigureEx{figures/rev-cor-zinc.pdf}{Deux H\textsuperscript{+} compensent la charge de Zn²\textsuperscript{+} ; les atomes et la charge totale sont conservés.}
\end{corrige}
\end{minipage}
\par\noindent\begin{minipage}{\linewidth}
\begin{corrige}{3}{Les réactifs}
\begin{questions}
\item Fer et ions H+.
\item Oui, si le contexte fournit une solution chlorhydrique et qu’il ne participe pas à la transformation simplifiée.
\end{questions}
\end{corrige}
\end{minipage}
\par\noindent\begin{minipage}{\linewidth}
\begin{corrige}{4}{Protection}
\begin{questions}
\item Elle constitue une barrière contre le milieu.
\item Elle peut exposer le métal à l’eau et au dioxygène.
\end{questions}
\pcFigureEx{figures/rev-cor-protection.pdf}{Une peinture protège en limitant le contact ; une rayure peut laisser agir le milieu sur le métal exposé.}
\end{corrige}
\end{minipage}
\par\noindent\begin{minipage}{\linewidth}
\begin{corrige}{5}{Le bon test}
\begin{questions}
\item Ion fer(II).
\item Il met en jeu un gaz inflammable et exige matériel, quantités et encadrement adaptés.
\end{questions}
\end{corrige}
\end{minipage}
\par\noindent\begin{minipage}{\linewidth}
\begin{corrige}{6}{Conserver la charge}
\begin{questions}
\item Deux charges positives de chaque côté.
\item Deux atomes H de chaque côté.
\end{questions}
\pcFigureEx{figures/rev-cor-zinc.pdf}{Deux H\textsuperscript{+} compensent la charge de Zn²\textsuperscript{+} ; les atomes et la charge totale sont conservés.}
\end{corrige}
\end{minipage}
\par\noindent\begin{minipage}{\linewidth}
\begin{corrige}{7}{Le pH évolue}
\begin{questions}
\item Elle diminue, mais la solution reste acide.
\item Oui, dans le contexte et le modèle étudiés.
\end{questions}
\pcFigureEx{figures/rev-cor-ph-acide.pdf}{Le pH augmente de 2 à 4, mais reste inférieur à 7 ; la solution est encore acide.}
\end{corrige}
\end{minipage}
\par\noindent\begin{minipage}{\linewidth}
\begin{corrige}{8}{Une pesée ouverte}
\begin{questions}
\item $0{,}06\unite g$.
\item Elles emporteraient aussi de la matière et empêcheraient d’attribuer toute la perte au gaz.
\end{questions}
\end{corrige}
\end{minipage}
\par\noindent\begin{minipage}{\linewidth}
\begin{corrige}{9}{Choisir un couple}
\begin{questions}
\item Dans ces conditions, l’humidité favorise la corrosion ; le contrôle des autres facteurs est nécessaire.
\item De l’eau désaérée maintenue à l’abri du dioxygène, comparée à B.
\end{questions}
\pcFigureEx{figures/rev-cor-corrosion.pdf}{Comparer des conditions contrôlées, avec des clous comparables et une même durée ; ne changer qu’un facteur dans la comparaison retenue.}
\end{corrige}
\end{minipage}
\par\noindent\begin{minipage}{\linewidth}
\begin{corrige}{10}{Ne pas tout changer}
\begin{questions}
\item Non, plusieurs causes changent simultanément.
\item Même métal, dimensions, durée et accès à l’air ; seule la présence de sel varie.
\end{questions}
\par\smallskip\noindent\textbf{Pour comprendre.}\enspace Choisir deux clous comparables, même volume d’eau, même accès à l’air, même température et même durée. Ne faire varier que la salinité. Comparer ensuite le même indicateur de corrosion ; une expérience isolée ne quantifie pas toutes les situations réelles.\par
\pcFigureEx{figures/rev-cor-corrosion.pdf}{Comparer des conditions contrôlées, avec des clous comparables et une même durée ; ne changer qu’un facteur dans la comparaison retenue.}
\end{corrige}
\end{minipage}
\par\noindent\begin{minipage}{\linewidth}
\begin{corrige}{11}{Un bilan de masse}
\begin{questions}
\item $8{,}0\unite g$.
\item La matière supplémentaire vient de l’environnement et doit être comptée dans les réactifs.
\end{questions}
\pcFigureEx{figures/rev-cor-masse-corrosion.pdf}{Une masse de matière issue du milieu s’ajoute au métal pour former les produits.}
\end{corrige}
\end{minipage}
\par\noindent\begin{minipage}{\linewidth}
\begin{corrige}{12}{Une conclusion trop large}
\begin{questions}
\item Non : l’essai ne concerne qu’un métal, un milieu et une durée.
\item Aucun dégagement visible n’a été observé pour ce cuivre dans ces conditions ; d’autres métaux ou acides peuvent se comporter autrement.
\end{questions}
\end{corrige}
\end{minipage}
\par\noindent\begin{minipage}{\linewidth}
\begin{corrige}{13}{Comparer des résultats}
\begin{questions}
\item Non : il faut connaître les conditions et la dispersion de mesure.
\item Répéter sur plusieurs clous comparables, avec témoins, même protocole et estimation des écarts.
\end{questions}
\end{corrige}
\end{minipage}
\par\noindent\begin{minipage}{\linewidth}
\begin{corrige}{14}{Choisir un modèle}
\begin{questions}
\item Non ; elle représente une réaction particulière du fer avec un acide.
\item La corrosion atmosphérique implique notamment eau et dioxygène et produit d’autres composés.
\end{questions}
\end{corrige}
\end{minipage}
\par\noindent\begin{minipage}{\linewidth}
\begin{corrige}{15}{Réparer plutôt que jeter}
\begin{questions}
\item Entretenir ou réparer une protection adaptée ; réemployer ou recycler le métal après tri.
\item Collecte, transport, séparation et remise en forme nécessitent de l’énergie.
\end{questions}
\end{corrige}
\end{minipage}
\par\noindent\begin{minipage}{\linewidth}
\phantomsection\label{cor-defi-5}\pcSolutionDefi{5}{Le portail près de la mer}\begin{questions}
\item Le revêtement isole le métal du milieu ; une rayure rompt cette barrière.
\item Comparer des pièces identiques, même durée et accès à l’air, dans de l’eau non salée ou salée ; répéter les essais.
\item Entretenir un revêtement adapté peut limiter la corrosion. Une expérience courte au laboratoire ne reproduit pas exactement les embruns, les variations de température et les contraintes réelles.
\end{questions}
\pcFigureEx{figures/rev-cor-protection.pdf}{Une peinture protège en limitant le contact ; une rayure peut laisser agir le milieu sur le métal exposé.}
\end{minipage}
\par\noindent\begin{minipage}{\linewidth}
\phantomsection\label{cor-activite-5}\pcSousTitre{Pistes pour l’investigation}
Air sec contre eau au contact de l’air teste le rôle de l’eau ; eau désaérée protégée contre eau au contact de l’air teste celui du dioxygène. L’eau salée doit être comparée à l’eau non salée avec mêmes pièces et mêmes durées. Pour un revêtement, seule la présence de protection varie. Des répétitions et une mesure de corrosion seraient nécessaires pour quantifier une accélération.\par\end{minipage}
